% =====================================================================
%  2bat-ud6-8.tex  —  SESSIÓ 8: És una diferència real o atzar?
%  (sense preàmbul: s'inclou des de main.tex)
% =====================================================================
\horatitol{Sessió 8 — És una diferència real o atzar?}%
{Activitat 3: davant d'una diferència entre dos grups, valorar de manera intuïtiva si és significativa o pot ser fruit de l'atzar.}

\subsection{Idea clau}

Treballem en parelles el cas de l'Activitat 3: es comparen els resultats de dos grups (per
exemple, l'eficàcia de dues intervencions socials) i s'observa una diferència. La pregunta
central: aquesta diferència és prou gran per considerar-la \textbf{real}, o podria ser fruit
de l'\textbf{atzar}? Aquesta és la lògica del \textbf{contrast d'hipòtesis}, que
introduïm de manera intuïtiva.

\begin{idea}
\textbf{Significativa o casual?} Davant d'una diferència observada entre dos grups, cal
preguntar-se sempre si és \textbf{significativa} (real) o \textbf{casual} (atzar del
mostreig). Tres claus per valorar-ho intuïtivament:
\begin{itemize}[leftmargin=1.4em, itemsep=2pt]
  \item \textbf{Mida de la diferència:} una diferència gran és més difícil d'atribuir a
        l'atzar que una de minúscula.
  \item \textbf{Mida de les mostres:} amb mostres grans, una diferència petita ja pot ser
        fiable; amb mostres petites, cal desconfiar.
  \item \textbf{Superposició d'intervals:} si els intervals de confiança dels dos grups
        \textbf{se superposen} molt, la diferència probablement \textbf{no} és
        significativa.
\end{itemize}
\textbf{La prudència clau:} no precipitar-se a dir «A és millor que B» davant de qualsevol
diferència.
\end{idea}

\subsection{Exemples}

\begin{exemple}[una diferència que probablement és casual]
Dues intervencions donen un èxit del $58\%\pm 5\%$ ($[53\%,63\%]$) i del $61\%\pm 5\%$
($[56\%,66\%]$). Els intervals \textbf{se superposen} molt (del $56\%$ al $63\%$), i la
diferència és de només $3$ punts. \textbf{Conclusió prudent:} no es pot afirmar que la
segona sigui millor; la diferència pot ser fruit de l'atzar.
\end{exemple}

\begin{exemple}[una diferència que sí sembla real]
Dos programes donen un èxit del $45\%\pm 4\%$ ($[41\%,49\%]$) i del $62\%\pm 4\%$
($[58\%,66\%]$). Els intervals \textbf{no} se superposen (hi ha un buit entre el $49\%$ i el
$58\%$) i la diferència és gran ($17$ punts). \textbf{Conclusió:} la diferència sembla
\textbf{real} (significativa): difícilment s'explica només per l'atzar.
\end{exemple}

\subsection{Exercicis resolts}

\begin{exresolt}[valorar amb els intervals]
El servei A té una satisfacció del $70\%\pm 3\%$ i el B del $74\%\pm 3\%$. La diferència és
significativa?
\solucio
Intervals: $A=[67\%,73\%]$ i $B=[71\%,75\%]$. \textbf{Se superposen} (del $71\%$ al $73\%$),
i la diferència és de només $4$ punts. \textbf{No} es pot afirmar que B sigui millor: la
diferència podria ser fruit de l'atzar del mostreig.
\resposta{No sembla significativa: els intervals se superposen.}
\end{exresolt}

\begin{exresolt}[el paper de la mida de la mostra]
Dues edicions d'un programa donen un $50\%$ i un $56\%$ d'èxit. En un cas les mostres són de
$30$ persones; en l'altre, de $2\,000$. En quin cas et fies més que la diferència sigui
real?
\solucio
Amb mostres de $30$ persones, el marge d'error és gran i una diferència de $6$ punts
fàcilment pot ser atzar. Amb mostres de $2\,000$, el marge és petit i la mateixa diferència
de $6$ punts és molt més fiable. Em fio més de la diferència amb les \textbf{mostres grans}
($2\,000$).
\resposta{Amb les mostres de $2\,000$: menys marge, diferència més fiable.}
\end{exresolt}

\subsection{Exercicis per fer}

\begin{exercicis}
\begin{exlist}
  \item Dos grups donen un $40\%\pm 4\%$ i un $43\%\pm 4\%$. Construeix els intervals i
        digues si se superposen. La diferència sembla significativa?
  \item Dos programes donen un $30\%\pm 3\%$ i un $48\%\pm 3\%$. Se superposen els
        intervals? La diferència sembla real?
  \item Una diferència de $2$ punts entre dos grups, amb mostres molt petites, et sembla
        fiable? Per què?
  \item Explica per què, amb mostres molt grans, fins i tot una diferència petita pot ser
        significativa.
  \item \textbf{(Prudència)} Un informe afirma «la intervenció A és millor que la B» a
        partir d'un $55\%$ enfront d'un $54\%$ amb mostres de $50$ persones. Estàs d'acord
        amb la conclusió? Com la matisaries?
  \item \textbf{(Raonament)} Per què davant de qualsevol diferència observada cal
        preguntar-se sempre si és real o casual, en lloc de concloure directament?
\end{exlist}
\end{exercicis}

% ---------------------------------------------------------------------
\fullsolucions{Sessió 8}
\begin{solucions}
\begin{sollist}
  \item Intervals $[36\%,44\%]$ i $[39\%,47\%]$: \textbf{se superposen} (del $39\%$ al
        $44\%$). La diferència de $3$ punts no sembla significativa (pot ser atzar).
  \item Intervals $[27\%,33\%]$ i $[45\%,51\%]$: \textbf{no} se superposen. La diferència
        de $18$ punts sembla \textbf{real} (significativa).
  \item No: amb mostres molt petites el marge d'error és gran, així que una diferència de
        només $2$ punts fàcilment pot ser fruit de l'atzar.
  \item Resposta oberta. Idea: amb mostres grans el marge d'error és petit i els intervals
        són estrets, de manera que una diferència petita ja pot deixar de superposar els
        intervals i esdevenir fiable (significativa).
  \item Resposta oberta. Idea: no s'hi hauria d'estar d'acord sense matisar: $1$ punt de
        diferència amb mostres de $50$ és perfectament atribuïble a l'atzar; caldria dir
        que la diferència no és concloent i que farien falta mostres més grans.
  \item Resposta oberta. Idea: perquè tota diferència observada en mostres pot venir en
        part de l'atzar del mostreig; concloure directament portaria a afirmar coses falses
        (per exemple, preferir una intervenció que en realitat no és millor).
\end{sollist}
\end{solucions}
